\chapter{项目介绍}\label{part:project}

“知微”可视分析系统旨在回答大语言模型应用中的一个核心工程问题：在不改动模型底层权重的前提下，提示词的组织方式、参考证据在上下文中的排布位置，以及执行反馈的信息粒度，如何共同决定任务最终的成败与计算代价。

\begin{figure}[!htbp]
\centering
\subfigure[五球时空交互入口]{\shot[0.48\linewidth]{\detokenize{figures/screenshots/gateway.jpg}}}\hfill
\subfigure[第一屏：研究问题与任务全景]{\shot[0.48\linewidth]{\detokenize{figures/screenshots/s1-default.jpg}}}
\caption{“知微”可视交互系统首屏概览}
\label{fig:ui-overview}
\end{figure}

如图~\ref{fig:ui-overview} 所示，系统采用深空星图与“景深”光学对焦作为视觉隐喻，将原本枯燥的 26 万余条离线日志封装为可探索的数字标本。用户既可以通过五个光晕入口直达“研究问题、实验设计、数据处理、核心发现、结论与反思”五个连续交互屏，也可随时切换至保留 44 张完整科学图稿的补充探索图册。

\paragraph{工程规模与统计口径}为清晰界定工作量，图~\ref{fig:project-stats} 展示了截至提交版本的代码库物理统计。全仓包含 \textbf{158 个核心源码文件}，合计 \textbf{37,534 行代码}（其中非空行 35,695 行），涵盖前端交互脚本（10,811 行）、专项实验分析（10,386 行）、离线清洗与构建管线（7,080 行）、样式系统（5,012 行）以及自动化测试套件（3,587 行）。

\begin{figure}[!htbp]
\centering
\includegraphics[max width=\linewidth]{\detokenize{figures/diagrams/d02-project-stats.pdf}}
\caption{项目数据规模与源码工作量分布}
\label{fig:project-stats}
\end{figure}

在数据口径方面，报告与系统严格区分三组物理计数，避免混淆统计分母：
\begin{itemize}
  \item \textbf{原五项实验统一表}（\texttt{data/processed/all\_records.csv.gz}）：共 \textbf{248,586 条}物理响应，包含推理提示、文档问答、代码生成、文档位置与历史 Mini 代码修复。
  \item \textbf{代码反馈补充批次}（\texttt{code-repair-extension-free-full-20261105-3rep-v2}）：独立保存的 \textbf{24,600 条}物理响应，覆盖 164 道题 $\times$ 2 模型 $\times$ 3 份共享初稿及六种反馈处理。
  \item \textbf{主展示页面选用口径}：新版“知微”主页面采用最新完整代码反馈补充批次替换历史 Mini 批次，共呈现 \textbf{267,610 条}物理响应（$248,586 - 5,576 + 24,600$），仓库全量累计物理响应为 \textbf{273,186 条}。
\end{itemize}

下面分别从项目背景与实验矩阵、研究需求与课程对照、四层系统架构与技术选型三个方面展开介绍。

\section{项目背景与实验设计}\label{sec:project-bg}

\subsection{受测模型与三类基准任务}

为了使实验结论兼具小参数稠密模型与大参数混合专家（MoE）模型的代表性，我们选取了四款主流开源及商业接口模型进行跨任务对照（详细公开规格及来源说明见附录~\ref{app:prompts}）。在本地推理与开源部署端，选用 \texttt{Qwen2.5-7B-Instruct}（7.61B 稠密）与 \texttt{Qwen3-8B}（8.2B 稠密，并在代码反馈补充实验中显式关闭额外思考模式）；在稀疏激活大模型端，选用 \texttt{DeepSeek-V4.1-Flash-1017}（552B 骨干，预填充激活 8B / 生成激活 16B）与 \texttt{HY4 preview}（770B 骨干，每 Token 激活 49B）。

\noindent\begin{minipage}[t]{0.52\linewidth}
\vspace{0pt}
任务设置覆盖课程要求的\textbf{三类典型 AI 交互场景}：第一类为数学推理与学科选择题，采用 \texttt{GSM8K}（小学算术应用题，200 题）、\texttt{MATH500}（竞赛数学题，200 题）与 \texttt{MMLU-Pro}（多学科高难选择题，200 题）；第二类为长上下文文档问答，基于 \texttt{NQ-Open} 构建 200 道问题与 10/20 篇干扰文档组合及闭卷基线；第三类为 Python 代码生成与多轮修复，分别基于 \texttt{CodePromptEval}（清洗审核后保留 194 题）与 \texttt{HumanEval+ Mini}（164 题）开展独立采样与多轮执行反馈对照。
\end{minipage}\hfill
\begin{minipage}[t]{0.45\linewidth}
\vspace{0pt}
\centering
\captionof{table}{受测模型与任务范围概览}\label{tab:models-tasks}
\ReportTableStyle
\begin{tabularx}{\linewidth}{@{}>{\raggedright\arraybackslash}p{1.9cm} >{\raggedright\arraybackslash}X@{}}
\toprule
\textbf{维度} & \textbf{实验范围} \\
\midrule
受测模型 & Qwen2.5 7B、Qwen3 8B、DeepSeek Flash、HY4 preview \\
推理/选择 & GSM8K、MATH500、MMLU-Pro（各 200 题） \\
文档问答 & NQ-Open 200 题（10/20 篇材料 + 闭卷） \\
代码任务 & CodePromptEval 194 题；HumanEval+ 164 题 \\
\bottomrule
\end{tabularx}
\end{minipage}

\subsection{31 种提示策略与五组实验矩阵}

在推理、文档问答与代码生成实验中，我们系统构造了覆盖\textbf{回答组织、语言模板、长度预算与表达格式}四大族系的 \textbf{31 种提示策略}（完整提示文本见附录~\ref{app:prompts} 表~\ref{tab:prompts-31}）。回答组织族包含只给答案（\texttt{NoCoT}）、默认分步（\texttt{DefaultCoT}）、简洁分步（\texttt{BeConcise}）以及步数限制（\texttt{StepLimit(1/2/5)}）；语言模板族包含中文作答（\texttt{ChineseCoT}）及其四档长度组合；长度预算族涵盖词数 7 档（1 至 100 词）、字母数 4 档（10 至 500 字母）和 Token 数 4 档（10 至 500 Token）；表达格式族则包含项目符号、缩写去元音、简化语法、移除空格与纯数字算式。

针对文档上下文位置效应与代码修复闭环机制，我们进一步增设了两项专题实验：其一为\textbf{文档证据位置专题}（3,600 条响应），固定 \texttt{NoCoT} 提示与 20 篇参考文档，在保持问题完全一致的前提下，将唯一包含答案的黄金文档依次置于第 1、3、5、8、10、13、15、18、20 位，精细刻画双模型的位置剖面曲线；其二为\textbf{代码执行反馈补充实验}（24,600 条响应），在共享同一份初始代码草稿的基础上，对比独立重采样、无反馈自行检查、仅返回通过状态、错误定位诊断、单个失败反例与多个失败反例六种处理臂在四轮修改中的演化轨迹。

\section{课程设计要求落实与多入口功能协同}\label{sec:requirements}

课程设计指导书对选项 B“可视交互系统”提出了六项明确的功能闭环要求：具备清晰的分析目标、打通从日志解析到行为分类的数据管道、实现不少于 3 个核心交互视图、提供至少一种跨视图有效联动、依托系统沉淀具体案例洞察，以及交付可独立运行的代码仓库与说明文档。我们在规划系统功能时，逐条核对了指导书第 5 至第 8 页的细则（完整评分点映射表见附录~\ref{app:validation}），并在此基础上做了更扎实的工程拓展。

\begin{figure}[!htbp]
\centering
\includegraphics[width=0.92\linewidth]{\detokenize{figures/diagrams/d02-project-stats.pdf}}
\caption{“知微”系统核心规模指标与跨任务评测覆盖概览}
\label{fig:project-stats}
\end{figure}

如图~\ref{fig:project-stats} 所示，系统在任务广度、视图丰富度与交互深度三个维度均大幅超出了基准线要求：

\paragraph{四类任务与多维行为分类管道}指导书要求至少设计 2～3 类任务且总消息数不少于 150 条。本项目实际覆盖了算术与竞赛数学（\texttt{GSM8K}、\texttt{MATH500}）、多学科专业选择（\texttt{MMLU-Pro}）、多文档开放问答（\texttt{NQ-Open}）以及 Python 函数生成与多轮测试修复（\texttt{CodePromptEval}、\texttt{HumanEval+}），样本规模达二十余万条。清洗脚本不仅提取了 Token 开销、响应耗时与停止原因，还通过规则分类器与沙箱执行器自动标定了 4 类外显输出组织形态、14 类细粒度代码运行错误以及多轮修复中的“初错修复 / 初对保持 / 过度修改回退”轨迹状态。

\paragraph{主叙事七大原生图版与工作台十三类交互图种}指导书第 5 页列举了时间线、柱状/饼图/堆叠图、错误恢复路径树、散点/热力图、分组箱线图、平行坐标图与原始文本详情面板等推荐图种。我们在主站（\texttt{index.html}）第四屏手写了七幅高信息密度的原生 SVG 交互图版，聚焦回答五项核心研究问题；同时在探索工作台（\texttt{explore.html}）完整实现了模型柱图、多维雷达图、31 策略热力图、规模开销气泡图、分组环形柱图、资料位置曲线、Plotly 分组箱线图、ECharts 平行坐标图、状态饼图与堆叠图、输出形态约束矩阵、HTTP 请求时间线、重试路径前缀树以及评分分歧桑基流图共计 \textbf{13 类交互视图}，让每一种推荐图种都承载真实的分析任务而非充当摆设。

\paragraph{五重跨视图交互联动机制}为了避免“点一下按钮只换一张静态图片”的假交互，我们在系统内打通了五条真实的数据联动链路：（1）第二屏点选任意提示词单元格，实时联动右侧实际英文指令模板、实验控制条件与样本覆盖统计；（2）第三屏点选 21 列数据字典字段，同步高亮上游原始表来源、清洗转换规则与真实匿名记录切片；（3）第四屏切换模型、基准题库或测试口径时，SVG 曲线、$R^2$ 对数拟合线、右上角配对比较卡的胜负平计数与 95\% 区间同步重绘；（4）点击“真实修改案例”滑出侧边抽屉，逐轮比对 JSON 测试反馈原文与前后代码行级 Diff；（5）在探索工作台中用鼠标框选（Brush）代码开销散点区域，右侧错误构成柱状图、可评样本计数与底部匿名记录导出集合立即按框选坐标过滤同步。


\subsection{面向课堂演示与自由钻取的双轨交互动线}

在完成全部视图组件开发后，团队结合日常排查数据的真实使用习惯（文档记录于 \texttt{docs/demo-guide.md} 与 \texttt{docs/product.md}），打磨了两条互补的分析操作动线，防止用户在丰富控件中迷失方向（图~\ref{fig:demo-paths}）：

\begin{figure}[!htbp]
\centering
\subfigure[主站第二屏：默认分步基线策略卡片状态]{\shot[0.485\linewidth]{\detokenize{figures/screenshots/s2-default.jpg}}}\hfill
\subfigure[探索台引导案例：输出长度直方分布与异常入口]{\shot[0.485\linewidth]{\detokenize{figures/screenshots/explore-demo.jpg}}}
\caption{面向结构化讲解的主站策略沙盘（左）与面向异常记录排查的探索台分布入口（右）}
\label{fig:demo-paths}
\end{figure}

\paragraph{动线一：五分钟主线验证动线（宏观规律 $\to$ 配对区间 $\to$ 单题代码 Diff）}适用于首次了解本项目的读者。从首页五球门户点击进入第一屏对比不同模型在 31 策略下的基准带宽；随后滑入第二屏（图~\ref{fig:demo-paths}a）交替点击 \texttt{DefaultCoT}、\texttt{BeConcise} 与 \texttt{CharLimit(10)} 核对原始英文指令；接着在第三屏查看 21 列统一表如何将 \texttt{T149} 环境排除项置空隔离；最后进入第四屏依次切换三大题库预算拟合线、9 位置剖面与代码修复轨迹，并在右侧配对卡中将“单个失败反例”对比“自行检查”，直接拉开侧滑抽屉检视 \texttt{HumanEval/54} 一行修好集合比较的真实代码差异。

\paragraph{动线二：探索式异常归因动线（分布框选 $\to$ 错误拆解 $\to$ 匿名子集导出）}适用于追查反常样本的研究者。从探索工作台（图~\ref{fig:demo-paths}b）切入“代码生成 $\to$ 输出开销”，用鼠标在 Vega-Lite 散点图上框选“高 Token 消耗但低 Pass@1”或“极短输出低通过率”的局部点簇，观察右侧 ECharts 柱状图是否由常规断言失败突变为高达八九成的 \texttt{extraction\_error}；随后点击对应错误柱打开底层匿名评分记录表，将命中条件导出为标准 CSV 文件做离线复核。两条动线通过顶栏导航按钮双向打通，确保叙事引导与自由探索无缝切换。

\section{四层零依赖工程架构与技术选型考量}\label{sec:architecture}

在动手搭建前端之前，我们小组认真讨论过一个很现实的问题：到底要不要上 React + Vite + Node 后端服务？讨论之后我们一致投了反对票。原因很简单：课程验收时老师和助教需要在不同电脑上快速拉起项目，如果强依赖特定版本的 Node.js、复杂的 \texttt{npm install} 编译链甚至在线 CDN 脚本，一旦遇到校园网波动或环境版本冲突，整个演示就会卡在安装依赖上。

基于这一工程判断，我们将全系统划分为职责边界清晰的\textbf{四层解耦架构}（图~\ref{fig:architecture}），实现了“离线重计算在 Python、在线轻交互在原生浏览器”的运行模式。

\begin{figure}[!htbp]
\centering
\includegraphics[width=0.95\linewidth]{\detokenize{figures/diagrams/d03-architecture.pdf}}
\caption{“知微”系统四层工程架构与离线快照流转关系}
\label{fig:architecture}
\end{figure}

\noindent
\begin{minipage}[c]{0.51\linewidth}
\setlength{\parindent}{2em}\setlength{\parfillskip}{0pt plus 0.78\linewidth}
\textbf{第一层（L1）为原始字节冻结与规范数据层}。所有采集落盘的 JSONL 原始响应与评测日志存放在 \texttt{data/raw/} 下并通过 Git LFS 管理；经脱敏与类型标准化后的分实验清洗表与 21 列总表保存在 \texttt{data/processed/}，由 \texttt{SHA256SUMS.txt} 锁死字节完整性。

\textbf{第二层（L2）为离线清洗、统计推断与科研绘图层}。分布在 \texttt{scripts/} 与 \texttt{experiments/} 下的 83 个 Python 脚本（合计 17,466 行代码，源码规模分布见图~\ref{fig:loc-arch}）承担全部耗时计算，包括 AST 代码规范化、沙箱隔离执行、5,000 次按题聚类 Bootstrap 抽样以及出版级矢量图导出。

\textbf{第三层（L3）为静态构建与防漂移快照层}。构建脚本 \texttt{build\_web\_app.py} 在校验源数据与图稿 SHA-256 指纹一致后，将前端所需的统计结果打包为 \texttt{app/data/} 下仅数兆大小的匿名 JSON 快照，并为静态资源注入 12 位内容哈希后缀，彻底消除浏览器缓存旧数据的隐患。
\end{minipage}\hfill
\begin{minipage}[c]{0.465\linewidth}
\centering
\includegraphics[width=\linewidth]{\detokenize{figures/charts/c08-loc-architecture.pdf}}
\captionof{figure}{全仓库 37,534 行源码分层构成统计}
\label{fig:loc-arch}
\end{minipage}

\medskip

\textbf{第四层（L4）为零外部网络依赖的交互前端层}。启动器 \texttt{app/launch.py} 仅调用 Python 3.10+ 标准库中的 \texttt{http.server} 与 \texttt{socket} 模块即可自动完成构建检核、端口绑定与本地服务拉起。在浏览器端，主站（\texttt{index.html}）完全不加载任何第三方图表框架，直接由 \texttt{zhiwei-gate.js} 与 \texttt{zhiwei.js} 操控 DOM 与原生 SVG \texttt{viewBox}；而探索工作台（\texttt{explore.html}）所使用的 ECharts、D3 7.9.0、Vega 5.30 / Vega-Lite 5.23 与 Plotly.js 3.1.1 共 2.45\,MB 库文件均以固定版本离线内置于 \texttt{app/assets/vendor/} 目录并附许可证校验。整套系统在完全断网的环境下依然能够毫秒级响应全部交互操作。

\subsection{离线压缩载荷封装与标准库无状态服务机制}

为了让包含近九万条代码评分与十五万条问答特征的探索工作台能够在不架设后台数据库的条件下流畅筛选，我们在 L3 构建层与 L4 浏览器运行时之间设计了\textbf{“Gzip 分片静态特征包 + 预编译全局运行时脚本”}的轻量传输方案：

\begin{itemize}
  \item \textbf{冷热数据分层加载}：首屏与主叙事五屏所需的汇总统计、置信区间与精选案例体积控制在数百 KB 以内（如 \texttt{research-charts.json} 仅 93\,KB、\texttt{code-feedback-extension.json} 仅 212\,KB），页面打开瞬间即可完成解析与首帧渲染；而体量较大的逐样本匿名特征明细（如 1.85\,MB 的 \texttt{features.json.gz} 与 1.18\,MB 的 \texttt{code-features.json.gz}）仅在用户切入探索工作台打开箱线图或记录抽屉时按需异步拉取解压。
  \item \textbf{零构建工具链的纯标准库启动器}：入口脚本 \texttt{app/launch.py} 不依赖 Flask、FastAPI 或任何第三方 Web 容器。启动时，它先调用内部构建器核验 \texttt{app/dist/} 快照指纹，随后通过标准库 \texttt{http.server.ThreadingHTTPServer} 在本地环回地址 \texttt{127.0.0.1} 上挂载多线程静态文件服务，自动配置 \texttt{.wasm}、\texttt{.woff2} 与 \texttt{.svg} 的标准 MIME 类型并处理端口冲突提示，真正做到克隆即运行。
\end{itemize}

